\documentclass[12pt,a4paper]{article}



\usepackage{polyglossia}
\setmainlanguage{french}

%\usepackage[T1]{fontenc}
%\usepackage[french]{babel}

\usepackage{amsmath, amssymb}

\usepackage{graphicx}
\usepackage{caption}

\usepackage{siunitx} % Pour les unités

\usepackage{geometry}
\geometry{top=2cm, bottom=2.5cm, left=3cm, right=3cm}
\setlength{\parindent}{0pt}

\usepackage{enumitem}

\usepackage{hyperref}

\usepackage{xcolor}

\usepackage{multicol}
\setlength{\columnseprule}{1pt}
\def\columnseprulecolor{\color{gray}}

\usepackage{tcolorbox}


\usepackage{import}
\usepackage{xifthen}
\usepackage{pdfpages}
\usepackage{transparent}
\usepackage{float}

\newcommand{\incfig}[1]{%
	\def\svgwidth{\columnwidth}
	\import{./figures/}{#1.pdf_tex}
}


\newcommand{\RomanNum}[1] %pour les chiffre romain
{\MakeUppercase{\romannumeral #1}}







\usepackage{tikz}
\usetikzlibrary{arrows.meta}
\usetikzlibrary{calc, angles, quotes, intersections}





\usepackage{cancel}


\usepackage{myunit}

\usepackage{stanli}




%commandes pour sigma
\newcommand{\sgma}[1]{\sigma_{#1}}
\newcommand{\sigmax}{\sigma_{\max}}
\newcommand{\sigmin}{\sigma_{\min}}

\newcommand{\sigmabar}[1]{\overline{\sigma}_{#1}}

%fonction pour tau
\newcommand{\taubar}[1]{\overline{\tau}_{#1}}









\title{Béton précontraint}
\author{}
\date{\today}

\begin{document}
	\counterwithin{section}{part} % pour que \section{} recommence de 1 aprés \part{}
	
	\maketitle
	
	\tableofcontents
	
	\newpage
	
	
	
	
	
	\part{Examen de rattrapage}
	
	
	
	\begin{center}
		
		
		\begin{tikzpicture}
			
			\point{A}{0}{0};
			\point{B}{6}{0};
			\point{C}{3}{0};
			\beam{2}{A}{B};
			
			\notation{1}{A}{A}[left];
			\notation{1}{B}{B}[right];
			
			\support{1}{A};
			\support{2oo}{B};
			
			\load{1}{C}[90][3];

			\foreach \xpos in {0,.5,...,6}
			\draw[->, thick] (\xpos,2) -- (\xpos,0);
			\draw (0,2) -- (6,2) node[left,above]{$q = 200 \cdot 10^{3} ~\unit{kg\per m}$};
			
			
			\draw (3,3.1) node[above]{\large F};
			
		\end{tikzpicture}
		
		
	\end{center}
	
	\begin{center}
		
		
		\begin{tikzpicture}
			
			\draw (0,0) rectangle (6,1.4);
			
			\draw[dashed] (-.5,.7) -- (6.5,.7);
			
			
			\draw[dashed,thin] (0,1) -- ++(0.5,0);
			\draw[dashed,thin] (6,.7) -- ++(.4,.3);
			
			\draw[thin] (0.4,.8) arc [start angle=-30, end angle=0, radius=5mm] node[right]{$22 \unit{\degree}$};
			\draw[thin] (6.3,.7) arc [start angle=-30, end angle=0, radius=5mm] node[right]{$21 \unit{\degree}$} ;
			
			\draw (0,1)  .. controls (1,.2) and (5,.2) .. (6,.7) ;
			
			\draw[|-|] (-.5,1) -- (-.5,.7) node[left,midway]{$\dfrac{h}{6}$};
			
			\draw[|-|] (3,.38) -- (3,.7) node[midway,right]{$0,6 ~\unit{m}$};
			
			\draw (8,0) rectangle (8.6,1.4);
			\draw[|-|] (7.7,0) -- (7.7,1.4) node[midway,left]{140};
			\draw[|-|] (8,-.3) -- (8.6,-.3) node[midway,below]{60};
			
		\end{tikzpicture}
	
		

	\end{center}




	
	\section{Évaluation des charges}
	
	\[ q = 200 \cdot 10^{3} ~\unit{kg\per m} = 2000 ~\unit{kN\per m} \]
	
	\[ \text{Poids propre de la poutre : } g = \rho_{ba} \cdot b \cdot h = 25 \cdot 0,6 \cdot 1,4 = 21 ~\unit{kN\per m} \]
	
	\[ F = 2500 ~\unit{kN} \]



	\section{Calcul des réactions et efforts}


	On va calculer pour chaque chargement séparément
	
	\subsection{Calcul avec le poids propre}
	
	\subsubsection{Réactions}
	
	\[ R_{H}^{A} = 0 \]
	
	{\centering par symétrie on aura : \par }
	
	\[ R_{V}^{A} = R_{V}^{B} = \dfrac{g \cdot l}{2} = \dfrac{21 \cdot 30}{2} = 315 ~\unit{kN} \]
	
	
	\subsubsection{Efforts tranchants}
	
	Poutre continue et uniformément charger donc on aura les efforts tranchants au appuie égale au réactions.
	
	\[ T_{g}(x = 0) =  - 315 ~\unit{kN} \]
	
	\[ T_{g}(x = 30 ~\unit{m}) = 315 ~\unit{kN} \]
	

	\vspace{\baselineskip}
	
	Depuis l'équation d'équilibre qu'on a eu a partir de la méthode des sections on a : $T_{g} = 21x - 315$
	
	
	\subsubsection{Moments}
	
	Charge uniformément répartie sur deux appuie simple, donc on aura le moment max en travée a x = 15m, les moments au appuie sont nul.
	
	\[ M_{g}(x = 15 ~\unit{m}) = \dfrac{g \cdot l^{2}}{8} = \dfrac{21 \cdot 30^{2}}{8} = 2362,5 ~\unit{kNm} \]
	
	\[ M_{g}(x = 0,x = 30 ~\unit{m}) = 0 \]
	
	\vspace{\baselineskip}
	
	L'équation : $M_{g} = 10,5x^{2} - 315x$
	
	\subsection{Calcul avec la surcharge}
	
	Chargement uniforme, donc calcul similaire que avec le poids propre.
	
	\subsubsection{Réactions}
	
	\[ R_{V}^{A} = R_{V}^{B} = \dfrac{q \cdot l}{2} = 3 \cdot 10^{4} ~\unit{kN} \]
	
	\subsubsection{Efforts tranchants}
	
	\[ T_{q}(x = 0) =  - 3 \cdot 10^{4} ~\unit{kN} \]
	
	\[ T_{q}(x = 30 ~\unit{m}) = 3 \cdot 10^{4} ~\unit{kN} \]
	
	\vspace{\baselineskip}
	
	L'équation : $T_{q} = 2000x - 30 \cdot 10^{4}$
	
	\subsubsection{Moments}
	
	
	\[ M_{q}(x = 15 ~\unit{m}) = \dfrac{q \cdot l^{2}}{8} = \dfrac{2000 \cdot 30^{2}}{8} = 225 \cdot 10^{3} ~\unit{kNm} \]
	
	\[ M_{q}(x = 0,x = 30 ~\unit{m}) = 0 \]
	
	\vspace{\baselineskip}
	
	L'équation : $M_{g} =  - 1000x^{2} + 3 \cdot 10^{4}x$


	\subsection{Calcul avec la force}
	

	\subsubsection{Réactions}

	
	\[ R_{v}^{A} = R_{v}^{B} = \dfrac{F}{2} = \dfrac{2500}{2} = 1250\kn \]

	
	\subsubsection{Effort tranchant}
	
	\begin{center}
		\begin{tikzpicture}
			
			\point{a}{0}{0};
			\point{0}{3}{0};
			
			\point{t}{3.7}{-.5};
			\point{m}{3.7}{.5};
			
			\beam{2}{a}{0};
			
			\support{1}{a};
			
			\load{1}{a}[-90];
			\load{1}{t}[90];
			\load{3}{m}[0][180];
			
			\dimensioning{1}{a}{0}{-2}[x];
			
			\node at (a) [below=1] {$R_{v}^{A}$};
			\node at (0) [right] {0};
			\node at (t) [below,right] {T};
			\node at (m) [above=.3] {M};
			
			\draw[dashed,thin] (3,-.5) -- (3,.5);
			
		\end{tikzpicture}
	\end{center}




	\[ \sum F/yy' = 0 \rightarrow R_{v}^{A} - T = 0 \]
	
	\[ T = 1250\kn \]
	
	\subsubsection{Moments}
	
	\[ \sum M/0 = 0 \rightarrow  - M + R_{v}^{A} \cdot x = 0 \]

	\[ M = R_{v}^{A} \cdot x = 1250 \cdot x = \begin{cases}
		(x = 0)  = 0 \\
		(x = L/2 = 15\meterl) = 18750\knxm
	\end{cases} \]
	
	
	\begin{tcolorbox}[colback=gray!20, colframe=black, title=Remarque]
		Par symétrie on aura la même chose de l'autre cote mais l'effort tranchant sera de signe inverse ($T =  - 1250\kn$)
	\end{tcolorbox}


	
	
	\subsection{La somme}
	
	\subsubsection{Réactions}
	
	\[ R_{v}^{A} = R_{vpoids}^{A} + R_{vsurcharge}^{A} + R_{vF}^{A} = 315+\num{3e4}+1250 \]
	
	\[ \fbox{$R_{v}^{A} = R_{v}^{B} \text{ (par symétrie) } = 31565\kn$} \]


	\subsubsection{Moments}
	
	
	\paragraph{Au point x = 0 et x = 30m (par symétrie)}

	\[ \fbox{$M = 0$} \]
	
	\paragraph{Au point x = 15m}
	
	\[ M = 2362,5 + \num{225e3} + 18750 \]

	\[ \fbox{$M' = 246112,5\knxm$} \]




	\section{Calcul de la précontrainte}
	
	
	Section : $S = 0,6 \cdot 1,4 = 0,84\msq$
	
	Inertie : $I = \dfrac{bh^{3}}{12} = \dfrac{0,6 \cdot 1,4^{3}}{12} = 0,1372\minertia$
	
	
	\subsection{Au point x = 0}
	
	
	\begin{center}
		\begin{tikzpicture}
			
			\draw[very thick] (0,0) -- (0,2) coordinate (cdg) -- (0,4) coordinate (p) -- (0,6);
			
			\draw[dash pattern=on 15pt off 5pt on 5pt off 5pt] (-3,2) -- (3,2) node[right] {C.D.G};
			\draw[dashed] (-3,4) coordinate (o) -- (3,4) ;
			
			\draw (0,4) .. controls (1,1) .. (4,0) node[right] {Câble de précontrainte};
			
			\draw (.3,4) arc [start angle=0 ,end angle=-80, radius=3mm] node[right] { $22\degree$};
			
			\draw[very thick, <-, green] (p) -- ++(120:1.5) coordinate (pp) node[above] {P};
			
			\draw[very thick,<-,green!50] (p) -- ++(0,1) node[right] {$P_{v}$};
			\draw[very thick,<-,green!50] (p) -- ++(-1,0) node[left] {$P_{h}$};
			
			
			\path[name path = P] (p) -- (pp);
			\path[name path = O] (p) -- (o);
			\path[name intersections = {of = P and O, by = I}];
			
			\pic [draw, "$22\degree$", angle eccentricity=1.5, angle radius=1cm]
			{angle = pp--I--o};
			
			\coordinate (O) at (0,2);
			\dimensioning{2}{O}{p}{-2.5}[$e = \dfrac{h}{6}$];
			
		\end{tikzpicture}
	\end{center}


	$P_{v}$ et $P_{h}$ sont la décomposition de P.
	
	\begin{tcolorbox}[colback=gray!20, colframe=black, title=Remarque]
		Pour pouvoir calculer les contrainte on doit avoir toutes les forces appliquer au CDG.
		
		Donc on doit déplacer la force $P_{h}$ vers le CDG et créer un moment qui n'est d'autre que $M_{P_{h}} = P_{h} \cdot e$
		
		
			\begin{center}
				\begin{tikzpicture}[scale=.6]
					
					\draw[very thick] (0,0) -- (0,2) coordinate (cdg) -- (0,4) coordinate (p) -- (0,6);
					
					\draw[dash pattern=on 15pt off 5pt on 5pt off 5pt] (-3,2) -- (3,2) node[right] {C.D.G};
					\draw[dashed] (-3,4) coordinate (o) -- (3,4) ;
					
					\draw (0,4) .. controls (1,1) .. (4,0) node[right] {Câble de précontrainte};
					
					\draw (.3,4) arc [start angle=0 ,end angle=-80, radius=3mm] node[right] { $22\degree$};
					
					\draw[very thick, <-, green] (p) -- ++(120:1.5) coordinate (pp) node[above] {P};
					
					\draw[very thick,<-,green!50] (p) -- ++(0,1.5) node[right] {$P_{v}$};
					\draw[very thick,<-,red!80] (cdg) -- ++(-1.5,0) node[left,below] {$P_{h}$};
					
					
					\path[name path = P] (p) -- (pp);
					\path[name path = O] (p) -- (o);
					\path[name intersections = {of = P and O, by = I}];
					
					\pic [draw, "$22\degree$", angle eccentricity=1.5, angle radius=1cm]
					{angle = pp--I--o};
					
					\coordinate (moment) at (-2,2);
					\node[above] at (moment) {$M_{P_{h}}$};
					\load{2}{moment};
					
					
					
				\end{tikzpicture}
			\end{center}
		
	\end{tcolorbox}



	\subsubsection{Calcul de P}
	
	\[ \sigmabar{P0} = \min \begin{cases}
		0,9 \cdot f_{peg} = 0,9 \cdot 1623 = 1460\mpa \\ \\
		0,8 \cdot f_{prg} = 0,8 \cdot 1814 = 1451\mpa
	\end{cases} \]
	
	\[ \sigmabar{P0} = 1451\mpa \]
	
	\[ \sigmabar{P0} = \dfrac{P}{S_{cable}} \Longrightarrow P = \sigmabar{P0} \cdot S_{cable} = 1451 \cdot 903 = 1310253\N \]
	
	\[ P = 1310\kn \]

	\subsubsection{Calcul des contraintes}
	
	\[ P_{h} = P \cdot \cos(22\degree) = 1310 \cdot \cos22\degree = 1214\kn \]
	
	\[ \sigma = \dfrac{N}{S} + \dfrac{M \cdot V}{I} \]
	
	\paragraph{Calcul de N} :
	
	\vspace{\baselineskip}
	
	Comme on a aucune force horizontale a part $P_{h}$ alors $N = P_{h} = 1214$
	 
	\paragraph{Calcul de M} :
	
	Calcul de $M_{P_{h}}$ moment due a la charge $P_{h}$ : 
	
	\[ M_{P_{h}} = P_{h} \cdot e = 1214 \cdot \dfrac{h}{6} = 1214 \cdot \dfrac{1,4}{6} = 283\knxm \]
	
	\vspace{\baselineskip}
	
	\[ M = M' + M_{P_{h}} = 0 + 283 = 283\knxm \]

	$M' = 0$ car les moments a x = 0 sont nul.
	
	\vspace{\baselineskip}
	
	\paragraph{Calcul des contraintes} :
	
	\[ \sgma{sup} = \dfrac{P_{h}}{S} + \dfrac{M \cdot \dfrac{h}{2}}{I} \]

	\[ \dfrac{h}{2} = 0,7\meterl \]

	\[ \sgma{sup} = \dfrac{1214}{0,84} + \dfrac{283 \cdot 0,7}{0,1372} = 2889,11\knms = 2,8\mpa \]
	
	\[ \sgma{inf} = \dfrac{1214}{0,84} + \dfrac{283 \times -(0,7)}{0,1372} = 1,36\knms = \num{1,36e-3}\mpa \]

	
	\subsubsection{Vérifications des contraintes}
	
	\[ \sigmabar{bt} = -\dfrac{f_{c28}}{11} \leq \sigma \leq \sigmabar{bc} = \dfrac{f_{c28}}{2} \]

	\[ -3,63\mpa \leq \sigma \leq 20\mpa \]

	Cette condition est vérifier pour $\sgma{sup}$ et $\sgma{inf}$.
	
	
	
	\subsection{Au point x = L/2 = 15m}
	
	
	Section : $S = 0,6 \cdot 1,4 = 0,84\msq$
	
	Inertie : $I = \dfrac{bh^{3}}{12} = \dfrac{0,6 \cdot 1,4^{3}}{12} = 0,1372\minertia$
	
	\[ P_{h} = P = 1310\kn \]
	
	\vspace{\baselineskip}
	
	\begin{center}
		\begin{tikzpicture}
			
			\draw (0,0) -- (7,0);
			\draw (0,3) -- (7,3);
			
			\draw[dash pattern=on 15pt off 5pt on 5pt off 5pt] (-3,1.5) -- (10,1.5) node[right] {C.D.G};
			
			\coordinate (cdg) at (3.5,1.5);
			
			\draw (-2,2) .. controls (-1,1) .. (1,.7) -- (6,.7) .. controls (8,1) .. (9,2);
			
			\coordinate (p) at (3.5,.7);
			
			\draw[very thick, red, <-] (p) -- ++(180:1.5) node[left,above] {P};
			
			\dimensioning{2}{cdg}{p}{4}[$e = 0,6$];
			
			
		\end{tikzpicture}
	\end{center}

	\vspace{\baselineskip}
	\[ \downarrow \text{On déplace la force P vers le CDG : } \]
	\vspace{\baselineskip}
	
	\begin{center}
		\begin{tikzpicture}
			
			\draw (0,0) -- (7,0);
			\draw (0,3) -- (7,3);
			
			\draw[dash pattern=on 15pt off 5pt on 5pt off 5pt] (-3,1.5) -- (10,1.5) node[right] {C.D.G};
			
			\coordinate (cdg) at (3.5,1.5);
			
			\draw (-2,2) .. controls (-1,1) .. (1,.7) -- (6,.7) .. controls (8,1) .. (9,2);
			
			\coordinate (p) at (3.5,.7);
			
			\draw[very thick, red, <-] (cdg) -- ++(180:1.5) node[left,above] {P};
			
			\load{3}{cdg};
			
			\node[below= .5 ] at (cdg) {$M_{P} = P \times e$};
			
			
		\end{tikzpicture}
	\end{center}

	
	
	
	\subsubsection{Calcul des contraintes}
	
	\[ \sigma = \dfrac{N}{S} + \dfrac{M \times V}{I} \]
	
	Avec :
	
	$N = P = 1310\kn$

	$M = M_{p} + M' = P \times e + 246112 = 1310 \times 0,6 + 246112 = 246898\knxm$
	
	\[ \sgma{sup}  = \dfrac{1310}{0,84} + \dfrac{246898 \times \dfrac{1,4}{2}}{0,1372} = 1261243\knms = 1261\mpa \]
	
	\[ \sgma{inf}  = \dfrac{1310}{0,84} + \dfrac{246898 \times -\dfrac{1,4}{2}}{0,1372} = -1258124\knms = -1258\mpa \]
	
	\subsubsection{Vérification des contraintes}
	
	\[ -3,63\mpa \leq \sigma \leq 20\mpa \]
	
	Les conditions ne sont pas vérifier il faut augmenter l'excentricité au max et si s'est toujours pas satisfait alors on augmente le nombre de câbles.
	
	
	
	
	
	\subsection{Au point x = 30m}
	
	Section : $S = 0,6 \cdot 1,4 = 0,84\msq$
	
	Inertie : $I = \dfrac{bh^{3}}{12} = \dfrac{0,6 \cdot 1,4^{3}}{12} = 0,1372\minertia$
	
	$P = 1310\kn$
	
	\vspace{\baselineskip}
	
	
	\begin{center}
		\begin{tikzpicture}[cm={-1,0,0,1,(0,0)}]
			
			\draw[very thick] (0,0) -- (0,2) coordinate (cdg) -- (0,4) coordinate (p) -- (0,6);
			
			\draw[dash pattern=on 15pt off 5pt on 5pt off 5pt] (-3,4) coordinate (o) -- (3,4)  node[right] {C.D.G};
			
			\draw (0,4) .. controls (1,1) .. (4,0) node[right] {Câble de précontrainte};
			
			\draw (.3,4) arc [start angle=0 ,end angle=-80, radius=3mm] node[right] { $21\degree$};
			
			\draw[very thick, <-, green] (p) -- ++(120:1.5) coordinate (pp) node[above] {P};
			
			\draw[very thick,<-,green!50] (p) -- ++(0,1) node[right] {$P_{v}$};
			\draw[very thick,<-,green!50] (p) -- ++(-1,0) node[left] {$P_{h}$};
			
			
			\path[name path = P] (p) -- (pp);
			\path[name path = O] (p) -- (o);
			\path[name intersections = {of = P and O, by = I}];
			
			\pic [draw, "$21\degree$", angle eccentricity=1.5, angle radius=1cm]
			{angle = o--I--pp};
			
			
		\end{tikzpicture}

	\end{center}
	
	
	
	\subsubsection{Calcul des contraintes}
	
	\[ \sigma = \dfrac{N}{S} + \dfrac{M \times V}{I} \]
	
	Avec :
	
	$N = P_{h} = P \times \cos21\degree = 1310 \times \cos21\degree = 1222\kn $
	
	$M = M' + M_{P_{h}} = 0 + 0$ car a x = 0 les moments sont nul et toutes les forces sont appliquer au CDG.
	
	
	\[ \sigma = \dfrac{N}{S} + \cancelto{0}{\dfrac{M \times V}{I}} \]
	
	
	\[ \sgma{sup} = \sgma{inf} = \dfrac{1222}{0,84} = 1454\knms = 1,4\mpa \]
	
	\subsubsection{Vérification des contraintes}
	
	\[ -3,63\mpa \leq \sigma \leq 20\mpa \]
	
	Conditions vérifier pour $\sgma{sup}$ et $\sgma{inf}$.
	
	
	
	
	
	
	
	
	
	
	
	
	
		
	
	
	
	
	
	
\end{document}